\documentclass[11pt,letterpaper]{article}
\usepackage[T1]{fontenc}
\usepackage{lmodern,microtype}
\usepackage[margin=1in]{geometry}
\usepackage{amsmath,amssymb,amsthm}
\usepackage{booktabs,tabularx}
\usepackage{xcolor}
\usepackage[colorlinks=true,linkcolor=blue!40!black,citecolor=blue!40!black,
            urlcolor=blue!40!black]{hyperref}

\newcommand{\manuscripttitle}{Minimum support of Boolean functions on a slice}
\newcommand{\manuscriptkeywords}{Boolean function, Johnson scheme, minimum support, inclusion matrix, uniform layer}
\hypersetup{
  pdftitle={Minimum support of Boolean functions on a slice},
  pdfauthor={Joseph M. Shunia},
  pdfsubject={Extremal Boolean families on uniform layers},
  pdfkeywords={\manuscriptkeywords}
}
\newtheorem{theorem}{Theorem}[section]
\newtheorem{lemma}[theorem]{Lemma}
\newtheorem{proposition}[theorem]{Proposition}
\newtheorem{corollary}[theorem]{Corollary}
\numberwithin{equation}{section}
\setlength{\emergencystretch}{1.5em}
\addtolength{\skip\footins}{6pt}

\title{\vspace{-1.5\baselineskip}%
  \manuscripttitle}
\author{Joseph M. Shunia\thanks{Independent researcher, Ann Arbor, Michigan, United States.
  Email: \href{mailto:jshunia@gmail.com}{\texttt{jshunia@gmail.com}}.}}
\date{November 22, 2026}

\begin{document}
\maketitle
\vspace{-\baselineskip}

\begin{abstract}
We determine the minimum support of a nonzero Boolean function of degree at most $d$ on a
uniform layer in several ranges of the conjecture of Kiermaier, Mannaert and Wassermann. We
prove their formula at degree two, at all widths at most six, and for every fixed degree on
sufficiently large ground sets. We also classify the minimizers one and two degrees below the
width. For fixed degree a density gap on the cube determines the candidate families. Near the
maximum degree an exact count on random matching cubes gives the extremizers and sharp
row-deletion thresholds for inclusion matrices.
\\[2mm]
{\bf 2020 Mathematics Subject Classification:} 05D05 (primary), 05E30, 06E30 (secondary).
\\[2mm]
{\bf Keywords:} \manuscriptkeywords.
\end{abstract}

\section{Introduction}\label{sec:intro}

Choose exactly five of ten objects. Requiring two specified objects to be chosen accepts 56 of
the 252 possibilities. Requiring three specified objects to be either all chosen or all
omitted accepts only 42. Both rules are quadratic. For the second rule the indicator in the
three membership bits is $1-x-y-z+xy+xz+yz$. Kiermaier, Mannaert and Wassermann calculated
this example and proposed an exact law for every degree and every layer \cite[Section 7]{kmw}.

Write $[n]=\{1,\ldots,n\}$. The $k$th layer $J(n,k)$ consists of its $k$-element subsets,
represented by membership bits $x_1,\ldots,x_n$. Its width is $\ell=\min(k,n-k)$. A family has
degree at most $d$ if a real polynomial of total degree at most $d$ agrees with its indicator
at every point of the layer. Let $M_{d}(n,k)$ be the minimum cardinality of a nonempty such
family. Replacing each set by its complement lets us assume $k\le n/2$. A $d$-pencil consists
of the sets containing a specified $d$-set.

Kiermaier, Mannaert and Wassermann's Conjecture~7 predicts the pencil size except on the
middle layer at even degree below the width. There the indicator that $d+1$ specified bits are
all equal has degree at most $d$, because its two highest-degree terms cancel. Define the
proposed minimum by
\begin{equation}\label{eq:minimum}
B_d(n,k)=
\begin{cases}
2\binom{2k-d-1}{k}, & n=2k,\ d\text{ even},\ d<k,\\[2pt]
\binom{n-d}{\ell-d}, & \text{otherwise}.
\end{cases}
\end{equation}
\begin{theorem}\label{thm:main}

Let $0\le d\le\ell$. Then $M_d(n,k)=B_d(n,k)$ in each of the following ranges:
\begin{enumerate}\setlength{\itemsep}{0pt}
\item $d=2$, or $\ell\le6$.
\item $\ell\ge B^d$ for an absolute constant $B>1$. For every fixed $d$ the equality also holds for all $k$ once $n\ge N(d)$.
\item $\ell-d\in\{1,2\}$. It also holds when $n=2k$, $k\ge3$ is odd and $d=k-3$.
\end{enumerate}
\end{theorem}

At fixed degree the proof starts with the junta theorem of Filmus and Ihringer \cite{fi}. We
identify the smallest cube supports and use the first correction to the middle-layer measure
to distinguish them. For fixed small widths we use the real inclusion-matrix theorem of Grosu,
Person and Szab\'o \cite{gps}. When $d$ is close to the width, the number of coordinates
involved is no longer bounded. Random matching restrictions also occur in the polynomial
distance lemma of Amireddy, Behera, Srinivasan and Sudan \cite{abss}. Here the inclusion
probability of each set depends only on its distance from a chosen support point. Keeping that
probability gives an exact counting identity and determines the extremizers.

\subsection{Equality for fixed degree}\label{sec:fixed}

Assume $k\le n/2$. In the large-width and large-ground-set ranges of
Theorem~\ref{thm:main}(2), the following list exhausts equality for $d\ge1$, up to
permutations of $[n]$.

\begin{center}\normalfont
\begin{tabularx}{\linewidth}{@{}llX@{}}
\toprule
Layer & Degree & Minimum families \\
\midrule
$n=2k$ & even & All of $d+1$ specified bits agree. \\
$n=2k+1$ & even & A $d$-pencil, or all of $d+1$ specified bits agree. \\
$n=2k$ & odd & A $d$-pencil or its complement image. Alternatively, one bit is prescribed and $d$ other specified bits all agree. \\
All other layers & either & A $d$-pencil. \\
\bottomrule
\end{tabularx}
\end{center}

For $d=1$ the condition that a single other bit agrees with itself is vacuous. For $d=0$ the
only nonempty family is the whole layer. For $d=\ell$ the minimum is one. The thresholds can
be enlarged so that $d=k$ occurs in the theorem only on the unbalanced fixed-width tail.

\begin{theorem}\label{thm:stability}

Fix $d\ge1$ and $\epsilon>0$. Along any sequence with $n\to\infty$ and $k/n\to1/2$, every
nonempty degree-at-most-$d$ family of density at most $(2-\epsilon)2^{-d}$ eventually has this
form. Fix $h$ bits, and on each of several disjoint odd blocks of size $r_j\ge3$ require one
of two complementary bit patterns. The fixed bits are disjoint from all blocks and
\[
h+\sum_j(r_j-1)=d.
\]

Conversely every such pattern has degree at most $d$ and density $(1+o(1))2^{-d}$ in this
limit. A pencil of degree $d-1$ has limiting density $2\cdot 2^{-d}$, so the factor two is
sharp.
\end{theorem}

The two thresholds in Theorem~\ref{thm:main}(2) allow the equality list above to omit finite
exceptions. For example, at $(n,k,d)=(8,4,3)$ one can independently complement members of a
pencil. Section~\ref{sec:high} describes exactly when this freedom occurs.

\section{Minimum support on a Boolean cube}\label{sec:cube}

\begin{lemma}\label{lem:cube}

On the uniform Boolean cube the density of a nonzero Boolean polynomial of degree at most $d$
is a positive integer multiple of $2^{-d}$. Equality at $2^{-d}$ holds precisely for the fixed
bits and complementary odd blocks described in Theorem~\ref{thm:stability}.
\end{lemma}

\begin{proof}

A Boolean function $F$ on a full cube has integer multilinear coefficients $c_S$, by iterated
finite differences at zero. Thus a nonzero function of degree at most $d$ has uniform density
$\sum_S c_S2^{-|S|}$ equal to a positive multiple of $2^{-d}$.

Suppose that the density is $2^{-d}$. Its degree is exactly $d$. Choose $d$ pivot variables
$x_1,\ldots,x_d$ whose monomial has coefficient $a\ne0$ and write $y$ for the other variables.
On every fixed-$y$ fiber the top pivot coefficient is $a$, so there is an accepted point in
each fiber. The total count gives exactly one per fiber. Write it
$b(y)=(b_1(y),\ldots,b_d(y))$. The top difference gives $a=(-1)^{d-\sum_i b_i(y)}$ on every
fiber. Consequently
\[
F(x,y)=a\prod_{i=1}^d(x_i+b_i(y)-1).
\]

Comparing the coefficient of the product of all pivots except $i$ gives $a(b_i(y)-1)$. It has
degree at most one in $y$. A Boolean affine function on the full cube is a constant, one bit,
or its complement. Two nonzero coefficients would give at least three distinct values on a
two-dimensional face.

Each pivot is therefore fixed or tied, with a possible bit flip, to one free variable. The
constant sign of $a$ forces each free variable that occurs to be used by an even number of
pivots. Such a variable and its pivots form an odd block of size at least three. This gives
exactly the form in Theorem~\ref{thm:stability}. Conversely the indicator of two complementary
point patterns on an odd block of size $r$ has degree at most $r-1$, since the leading terms
of its two point indicators cancel. Multiplication across disjoint blocks and fixed bits
proves the degree bound. Its density is exactly $2^{-d}$. There are
$L=d+\#\{\text{blocks}\}\le3d/2$ constrained bits.
\end{proof}

Multiplication by the parity character relates this cube equality problem to minimum-volume
trades \cite{trades}. The slice measure changes the relative sizes of the equality cases.

\subsection{Restricting the slice to its relevant coordinates}\label{sec:junta}

Filmus and Ihringer \cite[Theorem 1.1]{fi} give an absolute $C$ such that when $\ell\ge C^d$ every
Boolean function of degree at most $d$ on the slice depends on at most $J=d2^{d-1}$
coordinates. Enlarge $C$ to exceed one, and also require $\ell\ge J$. Pair each of the
relevant coordinates with a different outside coordinate. Assign the paired variables $x_i$
and $1-x_i$, and fix the remaining coordinates with the required number of ones. This embeds
the full relevant cube into the slice. Substituting into a degree-$d$ representative gives a
Boolean cube polynomial of degree at most $d$ which agrees with the family on every membership
pattern.

Let $u=n-2k$. Under a uniform random $k$-set, the probability of any prescribed pattern on $L$
coordinates with $w$ ones, divided by its fair cube probability, is
\[
D_{n,u}(L,w)=\frac{2^L(k)_w(n-k)_{L-w}}{(n)_L},
\]

where $(a)_j$ is a falling factorial. For fixed $d$ and $k/n\to1/2$ these ratios tend to one
uniformly for $L\le J$. The cube density gap and the characterization in
Section~\ref{sec:cube} now prove Theorem~\ref{thm:stability}.

For $u=0,1$, the logarithm of the ratio is the sum of $\log(1-(u+2j)/n)$ for $0\le j<w$, plus
the sum of $\log(1+(u-2j)/n)$ for $0\le j<L-w$, minus the sum of $\log(1-j/n)$ for $0\le j<L$.
In particular $|\log D|\le C_1(L+1)^2/n$ for $(L+1)^2/n$ sufficiently small, with an absolute
$C_1$. Choose $B$ so that $\ell\ge B^d$ meets the junta threshold and $J$ bound, and makes
every ratio with $L\le J$ lie between $0.9$ and $1.1$. This is possible since $(d2^{d-1}+1)^2$
is bounded by an absolute exponential in $d$. The pencil has cube density $2^{-d}$. Every cube
function of greater density has at least twice that much. Its slice mass therefore exceeds
that of the pencil. Every minimizer on these two layers belongs to Section~\ref{sec:cube}'s
list.

\subsection{The first correction on the two central layers}\label{sec:central}

For a pattern from Section~\ref{sec:cube} only $L\le3d/2$ bits enter. Expanding the same logarithms gives
\[
D_{n,u}(L,w)=1+\frac{L-(2w-L)^2-2u(2w-L)}{2n}
+O\!\left(\frac{(L+1)^4}{n^2}\right).
\]

The sum of the quadratic remainders in the logarithms is $O((L+1)^3/n^2)$. Exponentiating adds
$O((L+1)^4/n^2)$. All constants here are absolute and the bounds hold uniformly for $u=0,1$
when $(L+1)^2/n$ is small.

Condition on acceptance in the fair cube. The excess of ones over zeros among the constrained
bits is $2w-L=b_0+\sum_j b_j\varepsilon_j$, where $b_0$ is the fixed bits' excess, $b_j$ is
that of a chosen pattern on block $j$, and the $\varepsilon_j$ are independent fair signs.
Averaging shows that the slice density divided by $2^{-d}$ is
\[
1-\frac{\Phi_u}{2n}+O\!\left(\frac{(d+1)^4}{n^2}\right),
\qquad
\Phi_u=b_0^2+2u b_0-h+\sum_j(b_j^2-r_j).
\]

Two different integer scores are strictly ordered by this expansion for $n>C_2(d+1)^4$ with an
absolute $C_2$. Enlarge $B$ once more to ensure this. A larger score gives a smaller family.

Put $t_j=r_j-1$. The score is at most
\[
h(h-1+2u)+\sum_j t_j(t_j+1).
\]

Equality requires every free block to be homogeneous. All fixed bits must have the same value,
and for $u=1$ any fixed bits must be ones. Combining two nonempty free blocks' degree costs
increases the bound by $2t_it_j$. It therefore suffices to maximize
\[
(d-t)(d-t-1+2u)+t(t+1)
\]

over even $t$ from zero to $2\lfloor d/2\rfloor$. This strictly convex quadratic takes its
maximum only at endpoints. For even $d$, subtracting the value at zero from the value at $d$
gives $2d(1-u)$. For odd $d$, the difference between the values at $d-1$ and zero is
$-2u(d-1)$. These are precisely the types in Theorem~\ref{thm:main}(2).

For the exact ties at $n=2k$ and odd $d$, set complementation exchanges the two pencils. A
homogeneous $d$-block on its own has twice a pencil's size. Complementation pairs its members
with opposite values at any specified extra bit, so fixing that bit gives the exact pencil
size. At $n=2k+1$ and even $d$, the two parts of the homogeneous $(d+1)$-block have sizes
$\binom{2k-d}{k-d-1}$ and $\binom{2k-d}{k}$. By symmetry and Pascal's identity their sum is
$\binom{2k+1-d}{k-d}$. At $n=2k$ and even $d$ the two parts have equal size
$\binom{2k-d-1}{k}$. This proves both central cases, including the equality statements.

\section{Deletion of a coordinate}\label{sec:delete}

Write $x_D=\prod_{i\in D}x_i$. On the $k$th layer, each monomial $x_T$ of degree at most $d$
is a linear combination of degree-$d$ monomials, since
\[
\sum_{D\supseteq T,\ |D|=d}x_D=\binom{k-|T|}{d-|T|}x_T .
\]

The matrix of values of the $x_D$ has full real column rank if $d\le\min(k,n-k)$. Here is a
direct proof. In a vanishing linear combination, exchange two coordinates $i,j$ around any
$(k-1)$-set avoiding both. The differences $a_{T\cup\{i\}}-a_{T\cup\{j\}}$ give a vanishing
homogeneous combination of degree $d-1$ on the $(n-2,k-1)$ layer. The admissibility
inequalities persist, so induction from degree zero makes these differences zero. Every two
$d$-sets can be connected by such exchanges. The coefficients are equal, and evaluation
multiplies their common value by $\binom{k}{d}>0$.

Suppose $d<\ell$ and all accepted sets contain $i$. On the layer $x_i=0$ the coefficients of
$x_D$ with $i\notin D$ vanish by full column rank there. Factoring out $x_i$ gives a nonempty
Boolean family of degree at most $d-1$ on $(n-1,k-1)$. Complementing gives a corresponding
degree drop if all accepted sets omit $i$. The same argument for the containing case works
whenever $d\le\min(k,n-1-k)$ and $d\ge1$.

Alternatively every deletion $x_i=0$ is nonempty. Double-counting the omitted elements gives
\[
(n-k)|\mathcal F|=\sum_{i=1}^n|\mathcal F_{x_i=0}|
\ \ge\ n\,M_{d}(n-1,k).
\]

For $n\ge2k+2$, once the preceding layer has the pencil minimum, the right side divided by
$n-k$ is $\frac{n}{n-d}\binom{n-d}{k-d}$. This strictly exceeds a pencil. An extremizer must
contain a common element, and equality reduces to equality for the degree-$d-1$ family after
removing it.

Induct first on degree and then, for each fixed $k\ge B^d$, on $n$ starting at the two central
layers. The smaller-degree layer has $k-1\ge B^{d-1}$ when $d\ge2$ and $B\ge2$. In the
containing case its imbalance is at least three. Degree zero is immediate. Thus all off-center
minimizers in the stated range are pencils. The same induction, without a threshold, shows
that the numerical conjecture everywhere would follow from its middle and adjacent values for
all widths and degrees. It works as stated within any bounded range of widths as well. When
$d=k$, singleton families give the base directly.

For fixed $d$ there are finitely many remaining widths $\ell<B^d$. At fixed $k=\ell>d$,
Theorem 5 of Grosu, Person and Szabó \cite{gps} with rank deficiency one says that, for all sufficiently
large $n$, the fewest $k$-rows one can remove to lower the real column rank of the inclusion
matrix is $\binom{n-d}{k-d}$. Equality removes precisely the rows containing one $d$-set. The
rows outside the support of any nonzero slice polynomial annihilate its nonzero vector of
homogeneous coefficients. The rank theorem therefore supplies the fixed-width minimum and
uniqueness. Take the largest of the finitely many thresholds and enlarge it past $2B^d+2$.
Complementation completes Theorem~\ref{thm:main}(2).

\section{Quadratic families and layers of width at most six}\label{sec:small}

The support size of every Boolean family of degree at most $d$ is a multiple of
\[
\Delta(n,k,d)=\gcd_{0\le i\le d}\binom{n-i}{k-i}.
\]

We use this divisibility theorem and the elementary degree-one minimum $\binom{n-1}{\ell-1}$
for $\ell\ge1$. Both facts are given in \cite[Sections 2 and 7]{kmw}.

Start at $k=2$, where every singleton on the layer has degree two. Induct on $k\ge3$, treating
the middle and adjacent layer before using Section~\ref{sec:delete} for larger $n$. At $n=2k$
a family with a common element has at least the pencil size by the degree drop. Otherwise
double-counting the deletions gives at least $2M_{2}(2k-1,k-1)$ by complementation. The
preceding width gives $2\binom{2k-3}{k-3}=2\binom{2k-3}{k}$, exactly the homogeneous triple's
size.

Now set $n=2k+1$ and $P=\binom{2k-1}{k-2}$. A common element again gives at least $P$. Without
one, the middle-layer result gives the lower bound $A=\frac{2k+1}{k+1}\,2\binom{2k-3}{k}$.
Direct cancellation gives
\[
\frac{P-A}{P}=\frac{6}{(2k-1)(2k-2)},\qquad
\frac{\Delta(2k+1,k,2)}{P}=\frac{\gcd(2,k-1)}{k-1}.
\]

For the second identity the first three binomials defining $\Delta$, divided by $P/(k-1)$, are
$2(2k+1),2k,k-1$. The gap $P-A$ is strictly less than one unit of divisibility for every
$k\ge3$. Since $P$ is itself a multiple of $\Delta$, every family has size at least $P$.
Off-center propagation in Section~\ref{sec:delete} uses only the degree-one minimum at smaller
degree. This proves the numerical law for degree two without any size threshold.

It remains to determine the entries left open in Tables~1--3 of \cite{kmw}. The following
calculation uses the established entries of those tables.

At $(12,6,3)$ a common element gives at least $M_{2}(11,5)=84$. Otherwise the deletion average
gives at least $2M_{3}(11,5)=56$. Divisibility by 42 raises this bound to 84. At $(13,6,3)$ a
common element gives at least $M_{2}(12,5)=120$, and the alternative lower bound is
$13\cdot84/7=156$. The pencil attains 120.

At $(13,6,4)$ a common element gives at least $M_{3}(12,5)=36$. Otherwise all thirteen
deletions have size at least 14, and their sum is seven times the family size. The family size
is divisible by six. Thus the only possible size below 36 is 30. Every nonempty deletion has
size divisible by 14 on $(12,6)$. At size 30 its possibilities are 14 and 28. Some deletion
must have size 28 because $7\cdot30>13\cdot14$. Its complementary restriction would then have
size two on $(12,5)$, contradicting $M_{4}(12,5)=8$. A pencil attains 36.

The resulting values are as follows. For $n\ge2k+2$ the preceding pencil value and the
smaller-degree value give the entry directly by Section~\ref{sec:delete}. For the two degrees
where intermediate off-center rows are needed, the same recurrence gives $M_{2}(14,5)=220$ and
$M_{3}(14,5)=55$ from the preceding width-five rows.

\begin{center}\normalfont
\begin{tabularx}{\linewidth}{@{}lXl@{}}
\toprule
Degree & $(n,k)$ in the order shown & Minimum sizes \\
\midrule
2 & $(13,6),(14,6)$ & $330,\ 495$ \\
3 & $(13,5),(12,6),(13,6),(14,6),(15,6)$ & $45,\ 84,\ 120,\ 165,\ 220$ \\
4 & $(13,6),(14,6),(15,6)$ & $36,\ 45,\ 55$ \\
\bottomrule
\end{tabularx}
\end{center}

To prove the whole width-six assertion, take $k\le6$. The middle and adjacent base values for
$2\le d<k$ are supplied by Tables 1--3 and the preceding arguments, except for $k=6,d=5$. At
$(12,6,5)$ divisibility by seven gives the pencil value seven. At $(13,6,5)$ a common element
gives at least $M_{4}(12,5)=8$ and the deletion average otherwise gives at least
$13\cdot7/7=13$. A pencil has size eight. Degrees zero, one and $k$ give the boundary cases.
Section~\ref{sec:delete} propagates all these bases to arbitrary $n$ and finishes
Theorem~\ref{thm:main}(1).

\section{One and two degrees below the width}\label{sec:high}

Throughout the rest of the paper let $k\le m=n-k$ and put $t=k-d$. In this notation the pencil
size is $\binom{m+t}{t}$. At the even-degree middle-layer exception the minimum is
$2\binom{k+t-1}{t-1}$. On $J(16,8)$, for example, the minimum at degree seven is nine, although
each of the nine members of a seven-pencil may independently be replaced by its complement.

\begin{theorem}\label{thm:high}

For $t=1,2$ the minimum is \eqref{eq:minimum}. For $t=1$ its equality cases, up to permutation, are as follows.

\begin{itemize}
\item If $m=k$ and $k$ is odd, take one complementary pair.
\item If $m=k$ and $k$ is even, start with a $(k-1)$-pencil and independently complement any of its members.
\item If $m=k+1$ and $k$ is odd, take a pencil, or take a $k$-set $R$ together with all $k$-sets disjoint from $R$.
\item At all other parameters, take a pencil.
\end{itemize}

For $t=2$ the complete list is as follows.

\begin{itemize}
\item If $m=k$ and $k$ is even, require the $k-1$ elements of a fixed set to be all present or all absent.
\item If $m=k$ and $k$ is odd, take a pencil or its complement image. A further rule prescribes the membership of one element and requires $k-2$ other marked elements to be all present or all absent.
\item If $m=k+1$ and $k$ is even, take a pencil, or require $k-1$ marked elements to be all present or all absent.
\item At all other parameters, take a pencil.
\end{itemize}
\end{theorem}

Lists can coincide at degree zero or one. For even $k\ge4$, the number of distinct labelled
minimizers on $J(2k,k)$ at degree $k-1$ is $\binom{2k}{k-1}2^{k+1}$.

\subsection{The cost of a matching cube}\label{sec:weighted}

Choose $S$ in the support. A uniform injection from $S$ into its $m$-element complement pairs
each member with a distinct replacement. Switching any subset $A$ of the $k$ pairs gives a
vertex of a Boolean cube through $S$. Each original coordinate is constant, a switch bit, or
its complement, so polynomial degree never increases on restriction. Write $f(A)$ for the
value on the resulting cube and $w_j=\binom mj$.

A set $T$ at distance $j=|T\setminus S|$ is on this cube exactly when the images of
$S\setminus T$ form $T\setminus S$. These images form a uniform $j$-subset of the $m$ possible
replacements. Thus the probability of seeing $T$ is $1/w_j$. With $F$ the original support, we
have
\begin{equation}
|F|=\mathbb E\left[\sum_{A:\,f(A)\ne0}w_{|A|}\right].
\label{eq:weight}
\end{equation}

We call the sum the cost of the cube. Every injection has positive probability. If every cost
is at least $B$ and $|F|=B$, every cube through every $S\in F$ has cost $B$.

For degree at most $k-t$, the following alternating moments vanish for each $I\subseteq[k]$ with $|I|<t$:
\begin{equation}
\sum_{A\supseteq I}(-1)^{|A|}f(A)=0.
\label{eq:moments}
\end{equation}

Indeed, multiply the polynomial by the monomial prescribing that all bits in $I$ are one.
After multilinearization the degree is less than $k$. Every remaining monomial omits a bit,
and pairing the two values of that bit cancels its alternating sum. For $t=2$, a Boolean cube
support through the origin has at least four vertices. The zeroth moment gives equal numbers
of even and odd vertices. A support of size two is ruled out by a singleton moment at any bit
separating them.

We will also use the following observation. Suppose all nonzero vertices of a cube have
$|A|\le t$. In the retained-bit variables $y_i=1-1_{i\in A}$, the polynomial vanishes whenever
fewer than $d=k-t$ bits are one. Its multilinear coefficients of degrees less than $d$ vanish
by evaluating successively on each smaller subset of retained bits. It is therefore
homogeneous of degree $d$. Its coefficients of degree $d$ are its values at the vertices with
$|A|=t$. For a Boolean polynomial with $f(\varnothing)=1$, their sum is one and exactly one
equals one. Its support is the face obtained by freely switching $t$ specified pairs.

\subsection{Recovering a pencil from its restrictions}\label{sec:rigidity}

Suppose every cube through $S$ has support a face of $r$ free switches through the origin.
Form the zero-one matrix whose $(s,u)$ entry tells whether $S-s+u$ belongs to the support.
Each injection selects entries of sum $r$. If $m>k$, any row can use either of two columns
with the other $k-1$ rows held fixed. Its entries are consequently constant. The $r$ rows of
ones identify the free elements of $S$ for every injection. Every $T$ lies on some cube
through $S$, so the support is the pencil with core the other $k-r$ elements.

There is a balanced version for unordered pairs $\{T,T^c\}$ of $K$-sets on $2K$ points. Fix a
base $K$-set $S$. Suppose each orbit in the family has a representative within distance $r$ of
$S$, where $2r<K$, and on every perfect matching through $S$ these representatives form an
$r$-face. The nearby representative is unique, and a balanced matching cube contains a vertex
together with its complement. The same matrix of single switches now has constant sum on
permutations. Exchanging two assigned columns shows that every two rows differ by a constant
across columns. If a row contains both zero and one, all rows equal it. Otherwise each row is
constant. In the row case the near representatives form a pencil. In the column case they are
all the sets contained in $S$ together with $r$ specified outside elements, since only those
outside elements can be added. Complementing all representatives gives a pencil in this case
too. In particular the family of unordered pairs is the image of a full $(K-r)$-pencil.

The same proof applies to a balanced family whose elements are assigned an orientation
afterwards. It leaves the permissible orientations to be checked against the polynomial
degree.

\subsection{One degree below the width}\label{sec:one}

\begin{proof}[Proof of Theorem~\ref{thm:high} for $t=1$]

The zeroth moment forces a vertex at odd distance from the origin in every cube. For $m>k$ its
weight is at least $m$. For $m=k$ even, an odd distance lies strictly between $0$ and $k$, so
its weight is at least $k$. Equality in either bound on the total cost leaves just the origin
and a single odd vertex of weight $m$. Its distance is $1$ or $m-1$. These are the only
interior positions with $\binom mj=m$, with coincident positions counted once. For $m=k$ odd
the lower bound is two and the only second vertex of cost one is the antipode. Its mate is
$S^c$ for every matching. Complementary pairs attain degree at most $k-1$, since $\prod_{i\in
S}x_i+\prod_{i\in S}(1-x_i)$ loses its highest term for odd $k$.

When the only admissible distance is one, Section~\ref{sec:rigidity} gives an ordinary pencil
for $m>k$. In the even balanced case an extremizer contains at most one representative per
complementary pair. Indeed an equality cube through an accepted set has no second even vertex.
Every support orbit, oriented within distance one of $S$, now gives a single free switch on
each matching. Section~\ref{sec:rigidity} with $r=1$ identifies the orbits when $k\ge4$. Any
orientation works. Replacing $T$ with $T^c$ changes the indicator by the difference of
$\prod_{i\in T}(1-x_i)$ and $\prod_{i\in T}x_i$, of degree at most $k-1$ for even $k$. For
$k=2$ there are three complementary pairs, and a fixed one-element pencil meets each pair
once. The same description follows directly.

For even $k\ge4$ the core of the folded pencil is unique. Once one member of a lifted clique
of mutually distance-one sets is oriented, every other orientation in that clique is forced,
since distances one and $k-1$ differ. Hence the only two lifts which are such cliques are the
full pencil and its complement image. The latter consists of all $k$-subsets of a $(k+1)$-set
and has empty intersection. Only the former has a common $(k-1)$-core. Choosing the core and
then $k+1$ independent orientations gives the stated count.

For $m=k+1$ and $k$ odd, adjoin a new element $\infty$ to every member and view it as the
oriented representative of a complementary pair of $(k+1)$-sets on $2k+2$ elements. Any
matching at $S\cup\{\infty\}$ pairs $\infty$ with the one unused replacement of an injection
through $S$. On the original cube, a second vertex at distance one remains a neighbor. A
second vertex switching all $k$ original axes becomes, after complementation, the neighbor
switching only the $\infty$ axis. The near representatives thus give a one-dimensional face on
every balanced cube. For $k\ge3$, Section~\ref{sec:rigidity} gives a pencil on a core of size
$k$. Orient every orbit to contain $\infty$, then delete $\infty$. A core containing $\infty$
gives an ordinary pencil. A core $R$ avoiding $\infty$ gives $R$ and the $k$-sets disjoint
from it. Its indicator requires all or none of $R$, which has degree at most $k-1$ for odd
$k$. When $k=1$ the degree-zero function is constant. These constructions attain each lower
bound.
\end{proof}

\subsection{The possible squares two degrees below the width}\label{sec:two}

Suppose $k\ge4$. A pencil costs $B=\binom{m+2}{2}=1+2m+\binom m2$. For even $m=k$ the proposed
optimum is instead $E=2(k+1)$. We classify each cube at or below the relevant cost. In the
formulas below, $A$ is the set of switched axes and $[k]$ is the antipode.

\begin{center}\normalfont
\begin{tabularx}{\linewidth}{@{}p{.43\linewidth}X@{}}
\toprule
Parameters & Supports on the matching cube \\
\midrule
$m\ge k+2$, or $m=k+1$ with $k$ odd & $\varnothing,\{i\},\{j\},\{i,j\}$ \\
$m=k+1$, $k$ even & The preceding square, or $\varnothing,\{i\},[k]\setminus\{i\},[k]$ \\
$m=k$, $k$ even & $\varnothing,\{i\},[k]\setminus\{i\},[k]$ \\
$m=k$, $k$ odd & The ordinary square, or $\varnothing,\{i\},[k]\setminus\{j\},[k]\setminus\{i,j\}$ \\
\bottomrule
\end{tabularx}
\end{center}

Here $i\ne j$ whenever both indices occur. Each displayed support has exactly the proposed cost. It therefore suffices to prove the table.

\emph{Which distances can occur.} At budget $B$, a vertex with $3\le j\le m-3$ has cost at least $\binom m3>B$ for $m\ge8$. The inequality at eight persists since $m(m-1)(m-2)/[3(m+1)(m+2)]$ increases. For $m=5$ that interval is empty. For $m=6,7$ off the middle, the origin and two further vertices cost at least $1+2m$, to be added to $\binom m3$. For $m=k=7$ a weight of 35 together with the origin exhausts budget 36, whereas four vertices are required. Also at most one vertex of weight $\binom m2$, at distance $2$ or $m-2$, can occur when $m\ge5$. Two such weights and the origin already cost at least $B$ and a fourth vertex has positive weight.

For the even middle budget $E$, all distances $2,\dots,k-2$ are excluded. For $k\ge6$ their
weights are at least $\binom k2>E$. For $k=4$ a vertex of weight six, together with three
other vertices, costs at least $6+1+1+4>10$.

\emph{Off the middle.} If $m\ge k+3$, the support stays within distance two, and the homogeneous polynomial argument in Section~\ref{sec:weighted} gives a square. At $m=k+2$, the only additional possibility is the antipode, which costs $\binom m2$ and excludes every low pair. Write $e=(-1)^k$. If the antipode occurs, the singleton moment at $i$ gives $x_i=e$ for every singleton indicator $x_i$. The zeroth moment gives $\sum_i x_i=1+e$, implying $(k-1)e=1$, impossible for $k\ge4$.

At $m=k+1$, let $x_i$ indicate a singleton, $Y$ count low pairs, and $y_i$ count low pairs
containing $i$. Let $v_i$ indicate the vertex missing just $i$, $V=\sum v_i$, and let $z$
indicate the antipode. We have $Y+V\le1$. The two moment equations are
\[
1-\sum_i x_i+Y+e(z-V)=0,\qquad
x_i=y_i+e(z-V)+e v_i.
\]

If $Y=V=0$, these require $(k-1)ez=1$. If $Y=1$, the two indices in its pair and the indices
outside it force $z=0$, and precisely that square remains. If $V=1$, let $r$ be its missing
index. Then $x_r=ez$ and $x_i=e(z-1)$ elsewhere. For $e=1$ this forces $z=1$ and only $x_r=1$.
For $e=-1$ it forces $z=0$, whereupon the zeroth moment gives $k=3$. This proves both
off-middle rows.

\emph{At the middle.} For even $k$, only the origin, singletons, vertices missing one index, and the antipode can occur. With the same notation, the moments reduce to
\[
\sum_i x_i+V=1+z,\qquad x_i=z-V+v_i .
\]

Thus $(k-1)z-(k-2)V=1$, forcing $z=V=1$ and $x_i=v_i$. This gives the paired edge in the table.

For odd $k\ge5$, count also the high vertices missing a pair. Write $Q$ for their number and
$q_i$ for the number of missing pairs containing $i$. The bound on the two outer weights says
$Y+Q\le1$. Put $C=V-z-Q$. Directly summing \eqref{eq:moments} gives
\[
x_i+v_i=y_i+q_i+C,\qquad
\sum_i x_i=1+Y+C,\qquad (k-2)C=1+z-Y-Q .
\]

For the last identity, sum the first over $i$ and use $V=C+z+Q$. Its right side is in
$\{0,1,2\}$, so $C=0$, $z=0$ and $Y+Q=1$. If $Y=1$, then $V=0$ and the two singleton vertices
complete the ordinary square. If $Q=1$, then $V=1$ and exactly one low singleton occurs. Their
indices split the two positions in the missing pair, giving the last row. This completes the
cost proof.

\subsection{Equality two degrees below the width}\label{sec:eqhigh}

\begin{proof}[Completion of the proof of Theorem~\ref{thm:high}]

For $m>k$ in the ordinary-square case, the two free rows in Section~\ref{sec:rigidity} give
the pencil. For even $k$ and $m=k+1$, use the $\infty$ extension from Section~\ref{sec:one}.
The paired edge on the original cube folds to a square on axes $i$ and $\infty$. The ordinary
square already has two free axes. Every support orbit has a unique representative within
distance two of $S\cup\{\infty\}$ because $k+1\ge5$. Section~\ref{sec:rigidity} gives a
balanced pencil on $k-1$ elements. Orienting toward $\infty$ gives an ordinary $(k-2)$-pencil
if $\infty$ is in the core. If the core avoids $\infty$, a set is chosen precisely when it
contains or avoids the whole core. The latter indicator has degree at most $k-2$ since $k-1$
is odd. Its size is $\binom{k+2}{2}$ from the avoidance part plus $k+2$ from the containment
part, which sums to $\binom{k+3}{2}$.

For $m=k$ even, the table at each accepted root forces its complement to belong as well.
Folding each complementary pair into its representative at distance at most one leaves a
single free switch in every matching. By Section~\ref{sec:rigidity} the family consists
exactly of a $(k-1)$-pencil and the complements of all its members. This is the all-or-none
indicator on an odd core and has size $2(k+1)$.

For $m=k$ odd, apply the table at each accepted root. No complementary pair has both members
present. At a fixed root, the accepted sets project to nearby representatives at distance at
most two. In the last row the four nearby representatives are the empty switch, the two
singleton switches and their union. Thus every matching gives a square, and
Section~\ref{sec:rigidity} identifies the orbits as those of a $(k-2)$-pencil. It remains to
specify the orientations.

Let $R$ be its core, $U=[2k]\setminus R$, and write $h(ab)\in\{0,1\}$ for the decision to
complement the pencil member $R\cup\{a,b\}$. For four distinct elements of $U$, use two pairs
among them as free axes and match the remaining $k-2$ elements of $U$ bijectively to $R$. On
this cube the unswitched pencil occupies a square. Each complemented vertex changes its
contribution to the top alternating sum by minus twice its original sign, since $k$ is odd.
The top alternating sum of the oriented family vanishes. Using all pairings of the four
elements gives
\[
h(ab)+h(cd)=h(ac)+h(bd)=h(ad)+h(bc).
\]

Hence $h(iu)-h(ju)$ is independent of $u$ outside $i,j$. If every such difference is zero, $h$
is constant. Otherwise label $i,j$ so that the difference is one. Every edge from $i$ to
another vertex outside the pair has value one, and every such edge from $j$ has value zero.
For distinct $u,v$ outside the pair the same identity gives $h(uv)=1-h(ij)$. Thus the switched
edges are exactly a star or its complement. Constant decisions give the pencil and its
complement image. Switching the edges through $v$ requires $v$ to be absent and the elements
of $R$ to be all present or all absent. Complementing the whole family gives the rule with $v$
present. Since $|R|=k-2$ is odd, multiplying its all-or-none indicator by the prescribed bit
or its complement gives degree at most $k-2$. These orientations each choose exactly one
member of every orbit, giving $\binom{k+2}{2}$ sets. The off-middle pencils have the claimed
size by direct counting.

For $k=2$, $t=2$ means degree zero and the only nonempty family is the whole layer. For $k=3$,
$t=2$ means degree one. In a nonconstant linear Boolean function the difference of any two
coefficients is $0$ or $\pm1$, by exchanging those bits. Thus there are two coefficient levels
and the function is a constant plus the intersection size with one coefficient class. If both
classes have at least two elements, this intersection takes at least three consecutive values
because $k,m\ge2$. Consequently a nonconstant Boolean function prescribes one bit. Its two
orientations tie at the middle and prescribing presence wins on the smaller side. This
verifies the small cases of both the formula and the list.
\end{proof}

\begin{corollary}\label{cor:odd}

For every odd $k\ge3$, the minimum on $J(2k,k)$ at degree $k-3$ is $2\binom{k+2}{2}$.
\end{corollary}

\begin{proof}

A singleton support violates the zeroth moment on any matching through it. Two distinct
members differ on some coordinate, so both restrictions of the family at that coordinate are
nonempty. Each has degree at most $k-3$ on a layer of width $k-1$ in $2k-1$ points, after
complementing one restriction. Theorem~\ref{thm:high} at $t=2$ gives at least $\binom{k+2}{2}$
members in each restriction. Requiring $k-2$ marked elements to be all present or all absent
attains their sum, since $k-2$ is odd.
\end{proof}

\section{Polynomial support and the rank of inclusion matrices}\label{sec:rank}

\begin{theorem}\label{thm:rank}

Let $m>k$, $1\le t\le k$ and $d=k-t$. Suppose
\[
\binom mj\ge \binom{m+t}{t}\qquad(t<j\le k).
\]

Every nonzero complex polynomial function on $J(n,k)$ of degree at most $d$ has support at
least $\binom{m+t}{t}$. Equality holds precisely for nonzero scalar multiples of $d$-pencil
indicators. For a nontrivial abelian group $G$, a polynomial is a sum of square-free bit
monomials with coefficients in $G$, evaluated by adding the coefficients of the monomials
taking value one. The same minimum and equality list hold over $G$ if the displayed condition
holds and $2\binom mt\ge\binom{m+t}{t}$. In particular, for every such $G$ they hold whenever
$t\ge2$, $m\ge4t^2$ and $m-k\ge t$, including the boundary $m-k=t$. Each minimizing function
takes one fixed nonzero group value on a pencil and vanishes elsewhere.

In the complex range, and over every field in the group range, the minimum number of rows
whose deletion lowers the rank of the $k$-set versus $d$-set inclusion matrix is
$\binom{m+t}{t}$. Every minimum deletion removes a pencil.
\end{theorem}

For example, on $J(24,10)$ at degree seven, the first condition gives the exact minimum
$\binom{17}{3}=680$. The smallest binomial coefficient for $4\le j\le10$ is
$\binom{14}{4}=1001$, and $2\binom{14}{3}=728\ge680$. The example therefore works over every
coefficient group and every field in the rank conclusion.

\begin{proof}

Choose $S$ where the absolute value of the complex function is maximal, and divide by its
nonzero value there. Each cube has $f(\varnothing)=1$ and $|f(A)|\le1$. Write
$B_t=\binom{m+t}{t}$. If the cube has a support point with $t<|A|\le k$ under the first
hypothesis of Theorem~\ref{thm:rank}, that vertex and the origin together cost more than
$B_t$.

Otherwise the retained-bit polynomial is homogeneous of degree $d$, by
Section~\ref{sec:weighted}. Each degree-$d$ monomial survives on exactly $\binom tj$ vertices
with $j$ switched axes. Summing values on that level gives $\binom tj f(\varnothing)=\binom
tj$. The triangle inequality forces at least $\binom tj$ support points there. Vandermonde's
identity gives
\[
\sum_{j=0}^t\binom tj\binom mj=\binom{m+t}{t}.
\]

At equality the $j=t$ level contains just one nonzero value. It is the only nonzero
homogeneous coefficient, equal to one. Hence every equality cube is a face with constant value
on its support. Equation~\eqref{eq:weight} and the rectangular part of
Section~\ref{sec:rigidity} give one global pencil and the asserted constant value.

For $G$-valued functions under the additional inequality $2\binom mt\ge B_t$, root the cube at
any nonzero value. In a cube of cost at most $B_t$, the high-distance condition excludes all
distances greater than $t$. The homogeneous expansion in Section~\ref{sec:weighted} is valid
in $G$ by successive subtraction. Two nonzero coefficients would already cost $2\binom mt$ on
the $t$th level, in addition to the origin. There must be one nonzero coefficient because the
origin is nonzero. Thus each such cube is a single face with the same group value throughout
and has cost exactly $B_t$. Averaging and the row argument give the lower bound and all
equality cases.

To verify the uniform range, observe that $m\ge4t^2$ gives
\[
\frac{\binom mt}{B_t}=\prod_{i=1}^{t}\left(1-\frac{t}{m+i}\right)
\ge 1-\frac{t^2}{m+1}>\frac34 .
\]

It follows that $2\binom mt>B_t$ and $\binom m{t+1}>B_t$, the latter by multiplying the ratio
by $(m-t)/(t+1)$ for $t\ge2$. Unimodality and symmetry of binomial coefficients settle all
$t<j\le k$ if $m-k\ge t+1$. If $m-k=t$, only the antipode could escape that comparison. In a
cube costing at most $B_t$, a nonzero antipode excludes every point at distance $t$, since
both weights are $\binom mt$. Every distance strictly between $t$ and $k$ is already excluded.
Applying the group-valued cancellation \eqref{eq:moments} at each $(t-1)$-set $I$ now forces
$f(I)\ne0$, as only $I$ and the antipode can contribute. Their total cost across the distinct
vertices is at least
\[
\binom mt+\binom{k}{t-1}\binom m{t-1}>2\binom mt>B_t .
\]

Indeed $k=m-t$, $1\le t-1<k$, and $\binom{k}{t-1}\ge k>(k+1)/t=\binom mt/\binom m{t-1}$. This
rules out the antipode and completes the support theorem.

Finally, over any field every vector of the inclusion-matrix column space is a
degree-at-most-$d$ function, and each column itself is a pencil indicator. Deleting a row set
lowers rank exactly when a nonzero vector of the original column space vanishes on every
surviving row. The support lower bound and pencil equality list therefore give the stated
deletion result, including in positive characteristic.
\end{proof}

Theorem~\ref{thm:main} leaves open the general case in which both the degree and its distance
from the width grow. Even on the middle layer, the two moment conditions used for $t=2$ give
way to an increasing system of alternating equations. Determining which rooted cube supports
attain the least possible cost in that system would extend the matching argument to further
degrees.

\section*{AI Assistance Disclosure}

OpenAI's GPT-6.2 Astra was used in an autoresearch loop and reported this discovery to the author.

\begingroup\small
\begin{thebibliography}{9}
\bibitem{kmw}
M. Kiermaier, J. Mannaert, and A. Wassermann, \emph{The paired construction for Boolean
functions on the slice} (2025). \href{https://arxiv.org/abs/2510.02804}{arXiv:2510.02804}.

\bibitem{fi}
Y. Filmus and F. Ihringer, \emph{Boolean constant degree functions on the slice are juntas}
(2018). \href{https://arxiv.org/abs/1801.06338}{arXiv:1801.06338}.

\bibitem{gps}
C. Grosu, Y. Person, and T. Szab\'o, \emph{On the rank of higher inclusion matrices} (2013).
\href{https://arxiv.org/abs/1306.2197}{arXiv:1306.2197}.

\bibitem{abss}
P. Amireddy, A. R. Behera, S. Srinivasan, and M. Sudan, \emph{A near-optimal polynomial
distance lemma over Boolean slices}, in ICALP 2025, LIPIcs \textbf{334}, 11:1--11:17.
\href{https://doi.org/10.4230/LIPIcs.ICALP.2025.11}{doi:10.4230/LIPIcs.ICALP.2025.11}.

\bibitem{trades}
E. Ghorbani, S. Kamali, G. B. Khosrovshahi, and D. S. Krotov, \emph{On the volumes and affine
types of trades}. \href{https://arxiv.org/abs/1810.02296}{arXiv:1810.02296}.
\end{thebibliography}
\endgroup

\end{document}
